\SourceLeaf{232}allgemein, in dem Sinne dass $M$ die im innern Raume enthaltene Masse bedeutet, wohlverstanden, dass, wenn auch auf der Oberfläche selbst stetig vertheilte Massen sich befinden, diese den innern zugerechnet, oder davon ausgeschlossen werden müssen, jenachdem man für $\frac{\mathrm dV}{\mathrm dp}$ den auf die Aussenseite oder auf die Innenseite sich beziehenden Werth gewählt hat.

Sind demnach im Innern des Raumes gar keine Massen enthalten, so ist, wenn jedenfalls unter $\frac{\mathrm dV}{\mathrm dp}$ der auf die Innenseite sich beziehende Werth verstanden wird,
\[
\int\frac{\mathrm dV}{\mathrm dp}.\mathrm ds=0
\]
\Article{24}
Unter denselben Voraussetzungen, wie am Schluss des vorhergehenden Artikels, und indem wir den in Rede stehenden Raum mit $T$, und die in dem Elemente desselben $\mathrm dT$ durch die ausserhalb des Raumes oder auch nach der Stetigkeit in der Oberfläche vertheilten Massen entspringende ganze Kraft mit $q$ bezeichnen, haben wir folgenden wichtigen

\textsc{Lehrsatz.} Es ist
\[
\int V\frac{\mathrm dV}{\mathrm dp}.\mathrm ds=-\int qq\,\mathrm dT
\]
wenn das erste Integral über die ganze Fläche, das zweite durch den ganzen Raum $T$ ausgedehnt wird.

\textit{Beweis.} Indem wir rechtwinklige Coordinaten $x,y,z$ einführen, betrachten wir zuvörderst eine der Axe der $x$ parallele den Raum $T$ schneidende gerade Linie, wo also $y,z$ constante Werthe haben. Aus der identischen Gleichung
\[
\frac{\mathrm d}{\mathrm dx}(V\frac{\mathrm dV}{\mathrm dx})=(\frac{\mathrm dV}{\mathrm dx})^2+V\frac{\mathrm d\mathrm dV}{\mathrm dx^2}
\]
folgt, dass das Integral
\[
\int((\frac{\mathrm dV}{\mathrm dx})^2+V\frac{\mathrm d\mathrm dV}{\mathrm dx^2})\,\mathrm dx
\]
durch dasjenige Stück jener geraden Linie ausgedehnt, welches innerhalb $T$ fällt, der Differenz der beiden Werthe von $V\frac{\mathrm dV}{\mathrm dx}$ am Anfangs- und Endpunkte gleich wird, insofern die gerade Linie die Grenzfläche nur zweimal schneidet, oder allgemein $=\sum\varepsilon V\frac{\mathrm dV}{\mathrm dx}$, indem für $V\frac{\mathrm dV}{\mathrm dx}$ die einzelnen Werthe in den ver\SourceLeaf{233}schiedenen Durchschnittspunkten gesetzt werden, und $\varepsilon$ in den ungeraden Durchschnittspunkten (dem ersten, dritten u.\,s.\,f.) $=-1$, in den geraden $=+1$. Betrachten wir ferner längs dieser geraden Linie den prismatischen Raum, wovon das Rechteck $\mathrm dy.\mathrm dz$ ein Querschnitt, also $\mathrm dx.\mathrm dy.\mathrm dz$ ein Element ist, so wird das Integral
\[
\int((\frac{\mathrm dV}{\mathrm dx})^2+V\frac{\mathrm d\mathrm dV}{\mathrm dx^2})\,\mathrm dT
\]
ausgedehnt durch denjenigen Theil von $T$, welcher in jenen prismatischen Raum fällt, $=\sum\varepsilon V\frac{\mathrm dV}{\mathrm dx}.\mathrm dy.\mathrm dz$. Dieses Prisma scheidet aus der Grenzfläche zwei, oder allgemein eine gerade Anzahl von Stücken aus, und wenn jedes derselben mit $\mathrm ds$ bezeichnet wird, mit $\xi$ hingegen der Winkel zwischen der Axe der $x$ und der nach innen gerichteten Normale auf $\mathrm ds$, so ist $\mathrm dy.\mathrm dz=\pm\cos\xi.\mathrm ds$, das obere Zeichen für die ungeraden, das untere für die geraden Durchschnittspunkte genommen. Es wird folglich das obige Integral
\[
=-\sum V\frac{\mathrm dV}{\mathrm dx}.\cos\xi.\mathrm ds
\]
wo die Summation sich auf sämmtliche betreffende Flächenelemente bezieht. Wird nun der ganze Raum $T$ in lauter solche prismatische Elemente zerlegt, so werden auch die sämmtlichen correspondirenden Theile der Fläche diese ganz erschöpfen, und mithin
\[
\int((\frac{\mathrm dV}{\mathrm dx})^2+V\frac{\mathrm d\mathrm dV}{\mathrm dx^2})\,\mathrm dT=-\int V\frac{\mathrm dV}{\mathrm dx}.\cos\xi.\mathrm ds
\]
sein, indem die erste Integration durch den ganzen Raum $T$, die zweite über die ganze Fläche erstreckt wird. Offenbar ist nun $\cos\xi$ gleich dem partiellen Differentialquotienten $\frac{\mathrm dx}{\mathrm dp}$, indem $p$ die im Art.\ 23 festgelegte Bedeutung hat, und $x$ als Function von $p$ und zwei andern veränderlichen die einzelnen Punkte der Fläche von einander unterscheidenden Grössen betrachtet werden kann, folglich
\[
\int((\frac{\mathrm dV}{\mathrm dx})^2+V\frac{\mathrm d\mathrm dV}{\mathrm dx^2}).\mathrm dT=-\int V\frac{\mathrm dV}{\mathrm dx}.\frac{\mathrm dx}{\mathrm dp}.\mathrm ds
\]
Es ist übrigens von selbst klar, dass in dem Falle, wo die Fläche selbst Massen enthält, und also $\frac{\mathrm dV}{\mathrm dx}$ zwei verschiedene Werthe hat, hier immer der auf den innern Raum sich beziehende zu verstehen ist.

Durch ganz ähnliche Schlüsse findet man
\SourceLeaf{234}\[
\begin{gathered}
\int((\frac{\mathrm dV}{\mathrm dy})^2+V\frac{\mathrm d\mathrm dV}{\mathrm dy^2})\,\mathrm dT=-\int V\frac{\mathrm dV}{\mathrm dy}.\frac{\mathrm dy}{\mathrm dp}.\mathrm ds\\
\int((\frac{\mathrm dV}{\mathrm dz})^2+V\frac{\mathrm d\mathrm dV}{\mathrm dz^2})\,\mathrm dT=-\int V\frac{\mathrm dV}{\mathrm dz}.\frac{\mathrm dz}{\mathrm dp}.\mathrm ds
\end{gathered}
\]

Addirt man nun diese drei Gleichungen zusammen, und erwägt, dass im Raume $T$
\[
\begin{gathered}
\frac{\mathrm d\mathrm dV}{\mathrm dx^2}+\frac{\mathrm d\mathrm dV}{\mathrm dy^2}+\frac{\mathrm d\mathrm dV}{\mathrm dz^2}=0\\
(\frac{\mathrm dV}{\mathrm dx})^2+(\frac{\mathrm dV}{\mathrm dy})^2+(\frac{\mathrm dV}{\mathrm dz})^2=qq
\end{gathered}
\]
und an der Grenzfläche
\[
\frac{\mathrm dV}{\mathrm dx}.\frac{\mathrm dx}{\mathrm dp}+\frac{\mathrm dV}{\mathrm dy}.\frac{\mathrm dy}{\mathrm dp}+\frac{\mathrm dV}{\mathrm dz}.\frac{\mathrm dz}{\mathrm dp}=\frac{\mathrm dV}{\mathrm dp}
\]
so erhält man $\int qq\,\mathrm dT=-\int V.\frac{\mathrm dV}{\mathrm dp}.\mathrm ds$, welches unser Lehrsatz selbst ist, der unter Zuziehung des letzten Satzes des vorhergehenden Artikels noch allgemeiner sich so ausdrücken lässt
\[
\int qq\,\mathrm dT=\int(A-V)\frac{\mathrm dV}{\mathrm dp}.\mathrm ds
\]
wenn $A$ eine beliebige constante Grösse bedeutet.
\Article{25}
\textsc{Lehrsatz.} Wenn unter denselben Voraussetzungen, wie im vorhergehenden Artikel, das Potential $V$ in allen Punkten der Grenzfläche des Raumes $T$ einerlei Werth hat, so gilt dieser Werth auch für sämmtliche. Punkte des Raumes selbst, und es findet in dem ganzen Raume eine vollständige Destruction der Kräfte Statt.

\textit{Beweis.} Wenn in dem erweiterten Lehrsatze des vorhergehenden Artikels für $A$ der constante Grenzwerth des Potentials angenommen wird, so erhellt, dass $\int qq\,\mathrm dT=0$ wird, also nothwendig $q=0$ in jedem Punkte des Raumes $T$, mithin auch $\frac{\mathrm dV}{\mathrm dx}=0$, $\frac{\mathrm dV}{\mathrm dy}=0$, $\frac{\mathrm dV}{\mathrm dz}=0$, und folglich $V$ im ganzen Raume $T$ constant.
\Article{26}
\textsc{Lehrsatz.} Wenn von Massen, welche sich bloss innerhalb des endlichen Raumes $T$, oder auch, ganz oder theilweise nach der Stetigkeit vertheilt auf des\SourceLeaf{235}sen Oberfläche $S$ befinden, das Potential in allen Punkten von $S$ einen constanten Werth $=A$ hat, so wird das Potential in jedem Punkte $O$ des äussern unendlichen Raumes $T'$

\noindent erstlich, wenn $A=0$ ist, gleichfalls $=0$,

\noindent zweitens, wenn $A$ nicht $=0$ ist, kleiner als $A$ und mit demselben Zeichen wie $A$ behaftet sein.

\textit{Beweis.} I. Zuvörderst soll bewiesen werden, dass das Potential in $O$ keinen ausserhalb der Grenzen $0$ und $A$ fallenden Werth haben kann. Nehmen wir an, es finde in $O$ ein solcher Werth $B$ für das Potential Statt, und bezeichnen mit $C$ eine beliebige zugleich zwischen $B$ und $0$ und zwischen $B$ und $A$ fallende Grösse. Indem man von $O$ nach allen Richtungen gerade Linien ausgehen lässt, wird es auf jeder derselben einen Punkt $O'$ geben, in welchem das Potential $=C$ wird, und zwar so, dass die ganze Linie $OO'$ dem Raume $T'$ angehört. Dies folgt unmittelbar aus der Stetigkeit der Änderung des Potentials, welches, wenn die gerade Linie hinlänglich fortgesetzt wird, entweder von $B$ in $A$ übergeht, oder unendlich abnimmt, jenachdem die gerade Linie die Fläche $S$ trifft, oder nicht (vergl.\ die Bemerkung am Schlusse des 21.\ Artikels). Der Inbegriff aller Punkte $O'$ bildet dann eine geschlossene Fläche, und da das Potential in derselben constant $=C$ ist, so muss es nach dem Lehrsatze des vorhergehenden Artikels denselben Werth in allen Punkten des von dieser Fläche eingeschlossenen Raumes haben, da es doch in $O$ den von $C$ verschiedenen Werth $B$ hat. Die Voraussetzung führt also nothwendig auf einen Widerspruch.

Für den Fall $A=0$ ist hiedurch unser Lehrsatz vollständig bewiesen; für den zweiten Fall, wo $A$ nicht $=0$ ist, soweit, dass erhellt, das Potential könne in keinem Punkte von $T'$ grösser als $A$, oder mit entgegengesetztem Zeichen behaftet sein.

II. Um für den zweiten Fall unsern Beweis vollständig zu machen, beschreiben wir um $O$ als Mittelpunkt mit einem Halbmesser $R$, der kleiner ist als die kleinste Entfernung des Punkts $O$ von $S$, eine Kugelfläche, zerlegen sie in Elemente $\mathrm ds$, und bezeichnen das Potential in jedem Elemente mit $V$; das Potential in $O$ soll wieder mit $B$ bezeichnet werden. Nach dem Lehrsatze des 20.\ Artikels wird dann das über die ganze Kugelfläche ausgedehnte Integral
\[
\int V\mathrm ds=4\pi RRB,\quad\text{und folglich}\quad\int(V-B)\mathrm ds=0
\]
\SourceLeaf{236}Diese Gleichheit kann aber nur bestehen, wenn $V$ entweder in allen Punkten der Kugelfläche constant $=B$, oder wenn $V$ in verschiedenen Theilen der Kugelfläche in entgegengesetztem Sinne von $B$ verschieden ist. In der ersten Voraussetzung würde nach Art.\ 25 das Potential im ganzen innern Raume der Kugel und daher nach Art 21 im ganzen unendlichen Raume $T'$ constant, und zwar $=0$ sein müssen, im Widerspruche mit der Voraussetzung, dass es an der Grenze dieses Raumes, auf der Fläche $S$, von $0$ verschieden ist, und der Unmöglichkeit, dass es sich von da ab sprungsweise ändere. Die zweite Voraussetzung hingegen würde mit dem unter I: bewiesenen im Widerspruch stehen, wenn $B$ entweder $=0$ oder $=A$ wäre. Es muss daher nothwendig $B$ \textit{zwischen} $0$ und $A$ fallen.
\Article{27}
\textsc{Lehrsatz.} In dem Lehrsatze des vorhergehenden Artikels kann der erste Fall, oder der Werth $0$ des constanten Potentials $A$, nur dann Statt finden, wenn die Summe aller Massen selbst $=0$ ist, und der zweite nur dann, wenn diese Summe nicht $=0$ ist.

\textit{Beweis.} Es sei $\mathrm ds$ das Element der Oberfläche irgend einer den Raum $T$ einschliessenden Kugel, $R$ ihr Halbmesser, $M$ die Summe aller Massen und $V$ deren Potential in $\mathrm ds$. Da nach dem Lehrsatze des 20.\ Artikels das Integral $\int V\mathrm ds=4\pi RM$ wird, im ersten Falle oder für $A=0$ aber nach dem vorhergehenden Lehrsatze das Potential $V$ in allen Punkten der Kugelfläche $=0$ wird, im zweiten hingegen kleiner als $A$ und mit demselben Zeichen behaftet, so wird im ersten Fall $4\pi RM=0$, also $M=0$, im zweiten hingegen $4\pi RM$ und also auch $M$ mit demselben Zeichen behaftet sein müssen wie $A$. Zugleich erhellt, dass in diesem zweiten Falle $4\pi RM$ kleiner sein wird als $\int A\mathrm ds$ oder $4\pi RRA$, mithin $M$ kleiner als $RA$, oder $A$ grösser als $\frac{M}{R}$.

Der zweite Theil dieses Lehrsatzes, in Verbindung mit dem Lehrsatze des vorhergehenden Artikels, kann offenbar auch auf folgende Art ausgesprochen werden:

Wenn von Massen, die in einem von einer geschlossenen Fläche begrenzten Raume enthalten, oder auch theilweise in der Fläche selbst stetig vertheilt sind, die algebraische Summe $=0$ ist, und ihr Potential in allen Punkten der Fläche einen constanten Werth hat, so wird dieser Werth nothwendig selbst $=0$
\SourceLeaf{237}sein, zugleich für den ganzen unendlichen äussern Raum gelten, und folglich in diesem ganzen äussern Raume die Wirkung der Kräfte aus jenen Massen sich vollständig destruiren.
\Article{28}
Man wird sich leicht überzeugen, dass sämmtliche Schlüsse der beiden vorhergehenden Artikel ihre Gültigkeit behalten, wenn $S$ eine nicht geschlossene Fläche ist, und die Massen bloss in derselben enthalten sind. Hier fällt der Raum $T$ ganz weg; alle Punkte, die nicht in der Fläche selbst liegen, gehören dem unendlichen äussern Raume an, und wenn das Potential in der Fläche überall den constanten von $0$ verschiedenen Werth $A$ hat, wird es ausserhalb derselben überall einen kleinern Werth haben, der dasselbe Zeichen hat.

Das auf den ersten Fall, $A=0$, bezügliche bleibt zwar auch hier wahr, aber inhaltleer, da in diesem Fall das Potential $V$ in allen Punkten des Raumes $=0$ wird, mithin auch überall $\frac{\mathrm dV}{\mathrm dt}=0$, wenn $t$ irgend eine gerade Linie bedeutet, woraus man leicht nach Art.\ 18 schliesst, dass die Dichtigkeit in der Fläche überall $=0$ sein muss, also die Fläche gar keine Masse enthalten kann.

Diese letztere Bemerkung gilt übrigens allgemein, wenn die Massen bloss in der Fläche selbst enthalten sein sollen, auch wenn sie eine geschlossene ist, da offenbar nach dem Lehrsatz des 25.\ Artikels der Werth des Potentials in diesem Fall auch in dem ganzen innern Raume $=0$ sein wird.
\Article{29}
Ehe wir zu den folgenden Untersuchungen fortschreiten, in denen Massen, nach der Stetigkeit in eine Fläche vertheilt, eine Hauptrolle spielen, muss eine wesentliche bei der Vertheilung Statt findende Verschiedenheit hervorgehoben werden, indem nemlich entweder nur Massen von einerlei Zeichen (die wir der Kürze wegen immer als positiv betrachten werden) zugelassen werden, oder auch Massen von entgegengesetzten Zeichen. Ist eine Masse $M$ auf einer Fläche so vertheilt, dass auf jedes Element der Fläche $\mathrm ds$ die Masse $m\mathrm ds$ kommt, wo also nach unserm bisherigen Gebrauche $m$ die Dichtigkeit genannt, und $\int m\mathrm ds$ über die ganze Fläche ausgedehnt $=M$ wird, so nennen wir dies eine \textit{gleichartige} Vertheilung, wenn $m$ überall positiv, oder wenigstens nirgends negativ ist; wenn hingegen in einigen Stellen $m$ positiv, in andern negativ ist, so soll die Verthei\SourceLeaf{238}lung eine \textit{ungleichartige} Vertheilung heissen, wobei also $M$ nur die algebraische Summe der Massentheile, oder der absolute Unterschied der positiven und der negativen Massen ist. Ein ganz specieller Fall ungleichartiger Vertheilung ist der, wo $M=0$ wird, und wo es freilich anstössig scheinen mag, sich des Ausdrucks, die Masse $0$ sei über die Fläche vertheilt, noch zu bedienen.
\Article{30}
Es ist von selbst klar, dass, wie auch immer eine Masse $M$ über eine Fläche \textit{gleichartig} vertheilt sein möge, das daraus entspringende überall positive Potential $V$ in jedem Punkte der Fläche grösser sein wird, als $\frac{M}{r}$, wenn $r$ die grösste Entfernung zweier Punkte der Fläche von einander bedeutet: diesen Werth selbst könnte das Potential nur in einem Endpunkte der Linie $r$ haben, wenn die ganze Masse in dem andern Endpunkte concentrirt wäre, ein Fall, der hier gar nicht in Frage kommt, indem nur von stetiger Vertheilung die Rede sein soll, wo jedem Elemente der Fläche $\mathrm ds$ nur eine unendlich kleine Masse $m\mathrm ds$ entspricht. Das Integral $\int Vm\mathrm ds$ über die ganze Fläche ausgedehnt, ist also jedenfalls grösser als $\int\frac{M}{r}m\mathrm ds$ oder $\frac{MM}{r}$, und so muss es nothwendig eine gleichartige Vertheilungsart geben, für welche jenes Integral einen Minimumwerth hat. Es mag nun hier im Voraus als eines der Ziele der folgenden Untersuchungen bezeichnet werden, zu beweisen, dass bei einer solchen Vertheilung, wo $\int Vm\mathrm ds$ seinen Minimumwerth erhält, das Potential $V$ in jedem Punkte der Fläche einerlei Werth haben wird, dass dabei keine Theile der Fläche leer bleiben können, und dass es nur eine einzige solche Vertheilung gibt. Der Kürze wegen wollen wir aber die Untersuchung schon von Anfang an in einer weiter umfassenden Gestalt ausführen.
\Article{31}
Es bedeute $U$ eine Grösse, die in jedem Punkte der Fläche einen bestimmten endlichen nach der Stetigkeit sich ändernden Werth hat. Es wird dann das Integral
\[
\Omega=\int(V-2U)m\mathrm ds
\]
über die ganze Fläche ausgedehnt, zwar nach Verschiedenheit der gleichartigen Vertheilung der Masse $M$, sehr ungleiche Werthe haben können; allein offenbar
\SourceLeaf{239}muss für Eine solche Vertheilungsart ein Minimumwerth dieses Integrals Statt finden. Es soll nun ein Beweis gegeben werden für den

\textsc{Lehrsatz,} dass bei solcher Vertheilungsart

1. die Differenz $V-U=W$ überall in der Fläche, wo sie wirklich mit Theilen von $M$ belegt ist, einen constanten Werth haben wird;

2. dass, falls Theile der Fläche dabei unbelegt bleiben, $W$ in denselben grösser sein muss, oder wenigstens nicht kleiner sein kann, als jener constante Werth.

I. Zuvörderst soll bewiesen werden, dass wenn anstatt einer Verthei\-lungsweise eine andere unendlich wenig davon verschiedene angenommen wird, indem $m+\mu$ an die Stelle von $m$ gesetzt wird, die daraus entspringende Variation von $\Omega$ durch $2\int W\mu\mathrm ds$ ausgedrückt werden wird.

In der That ist, wenn wir die Variationen von $\Omega$ und $V$ mit $\delta\Omega$ und $\delta V$ bezeichnen,
\[
\delta\Omega=\int\delta V.m\mathrm ds+\int(V-2U)\mu\mathrm ds
\]
Allein zugleich ist $\int\delta V.m\mathrm ds=\int V\mu\mathrm ds$, wie leicht aus dem Lehrsatze des 19.\ Artikels erhellt, indem $\delta V$ nichts anders ist, als das Potential derjenigen Massenvertheilung, wobei $\mu$ die Dichtigkeit in jedem Flächenelemente vorstellt, und also was hier $V,m,\delta V,\mu$ ist, dort für $V,K,v,k$ angenommen werden kann, so wie $\mathrm ds$ zugleich für $\mathrm dS$ und $\mathrm ds$. Es wird folglich
\[
\delta\Omega=\int(2V-2U)\mu\mathrm ds=2\int W\mu\mathrm ds
\]

II. Offenbar sind die Variationen $\mu$ allgemein an die Bedingung geknüpft, dass $\int\mu\mathrm ds=0$ werden muss; für die gegenwärtige Untersuchung aber auch noch an die zweite, dass $\mu$ in den unbelegten Theilen der Fläche, wenn solche vorhanden sind, nicht negativ sein darf, weil sonst die Vertheilung aufhören würde, eine gleichartige zu sein.

III. Nehmen wir nun an, dass bei einer bestimmten Vertheilung von $M$ ungleiche Werthe der Grösse $W$ in den verschiedenen Theilen der Fläche Statt finden. Es sei $A$ eine Grösse, die zwischen den ungleichen Werthen von $W$ liegt; $P$ das Stück der Fläche, wo die Werthe von $W$ grösser, $Q$ dasjenige, wo sie kleiner sind, als $A$; es seien ferner $p,q$ gleich grosse Stücke der Fläche, jenes zu $P$, dieses zu $Q$ gehörig. Dies vorausgesetzt, legen wir der Variation
\SourceLeaf{240}von $m$ überall in $p$ den constanten negativen Werth $\mu=-\nu$, in $q$ hingegen überall den positiven $\mu=\nu$, und in allen übrigen Theilen der Fläche den Werth $0$ bei. Offenbar wird hiedurch der ersten Bedingung in II Genüge geleistet; die zweite hingegen wird noch erfordern, dass $p$ keine unbelegte Theile enthalte, was immer bewirkt werden kann, wenn nur nicht das ganze Stück $P$ unbelegt ist.

Der Erfolg hievon wird aber sein, dass $\delta\Omega$ einen negativen Werth erhält, wie man leicht sieht, wenn man diese Variation in die Form $2\int(W-A)\mu\mathrm ds$ setzt.

Es erhellt hieraus, dass wenn bei einer gegebenen Vertheilung entweder in dem belegten Stücke der Fläche ungleiche Werthe von $W$ vorkommen, oder wenn, bei Statt findender Gleichheit der Werthe in dem belegten Stücke, kleinere in dem nichtbelegten Theile angetroffen werden, durch eine abgeänderte Vertheilung eine Verminderung von $\Omega$ erreicht werden kann, und dass folglich bei dem Minimumwerthe nothwendig die in obigem Lehrsatze ausgesprochenen Bedingungen erfüllt sein müssen.
\Article{32}
Wenn wir jetzt für unsern speciellen Fall (Art.\ 30), wo $U=0$ ist, also $W$ das blosse Potential der auf die Fläche vertheilten Masse, und $\Omega$ das Integral $\int Vm\mathrm ds$ bedeutet, mit dem Lehrsatze des vorhergehenden Artikels den im 28.\ Artikel angeführten verbinden, so folgt von selbst, dass bei dem Minimumwerth von $\int Vm\mathrm ds$ die Fläche gar keine unbelegte Theile haben kann; denn sonst würde, auch wenn die ganze Fläche eine geschlossene ist, der belegte Theil eine ungeschlossene und hinsichtlich derselben der unbelegte Theil als dem äussern Raume angehörig zu betrachten sein, mithin darin nach Art.\ 28 das Potential einen kleinern Werth haben müssen als in der belegten Fläche, während der Lehrsatz des vorhergehenden Artikels einen kleinern Werth ausschliesst.

Es ist also erwiesen, dass es eine gleichartige Vertheilung einer gegebenen Masse über die ganze Fläche gibt, wobei kein Theil leer bleibt, und woraus ein in allen Punkten der Fläche gleiches Potential hervorgeht. Was zum vollständigen Beweise des im 30.\ Artikel aufgestellten Lehrsatzes jetzt noch fehlt, nemlich die Nachweisung, dass es nur Eine dies leistende Vertheilungsart geben kann, wird weiter unten als Theil eines allgemeineren Lehrsatzes erscheinen.

Dass, wenn der Minimumwerth für $\int Vm\mathrm ds$ Statt finden soll, kein Theil der Fläche unbelegt bleiben darf, kann offenbar auch so ausgedrückt werden:
\SourceLeaf{241}Bei jeder Vertheilung, wobei ein endliches Stück der Fläche leer bleibt, erhält das Integral $\int Vm\mathrm ds$ einen Werth, der den Minimumwerth um eine endliche Differenz übertrifft.
\Article{33}
Der eigentliche Hauptnerv der im 31.\ Artikel entwickelten Beweisführung beruht auf der Evidenz, mit welcher die Existenz eines Minimumwerths für $\Omega$ unmittelbar erkannt wird, solange man sich auf die gleichartigen Vertheilungen einer gegebenen Masse beschränkt. Fände eine gleiche Evidenz auch ohne diese Beschränkung Statt, so würden die dortigen Schlüsse ohne weiteres zu dem Resultate führen, \textit{dass es allemal, wenn nicht eine gleichartige, doch eine ungleichartige Vertheilung der gegebenen Masse gibt, für welche $W=V-U$ in allen Punkten der Fläche einen constanten Werth erhält}, indem dann die zweite Bedingung (Art.\ 31.\ II) wegfällt. Allein da jene Evidenz verloren geht, sobald wir die Beschränkung auf gleichartige Vertheilungen fallen lassen, so sind wir genöthigt, den strengen Beweis jenes wichtigsten Satzes unserer ganzen Untersuchung auf einem etwas künstlichern Wege zu suchen. Der folgende scheint am einfachsten zum Ziele zu führen.

Wir betrachten zunächst drei verschiedene Massenvertheilungen, bei welchen wir anstatt der unbestimmten Zeichen für Dichtigkeit $m$ und Potential $V$ folgende besondere gebrauchen:
\[
\begin{array}{rll}
\mathrm{I.}&m=m^0,&V=V^0\\
\mathrm{II.}&m=m',&V=V'\\
\mathrm{III.}&m=\mu,&V=v
\end{array}
\]
Die Vertheilung I ist diejenige gleichartige der positiven Masse $M$, für welche $\int Vm\mathrm ds$ seinen Minimumwerth erhält.

II. ist die gleichartige Vertheilung derselben Masse $M$, für welche $\int(V-2\varepsilon U)m\mathrm ds$ seinen Minimumwerth erhält, wo $\varepsilon$ einen beliebigen constanten Coëfficienten bedeutet.

III. hängt so von I und II ab, dass $\mu=\frac{m'-m^0}{\varepsilon}$, und ist alo eine ungleichartige Vertheilung, in welcher die Gesammtmasse $=0$ wird.

Es ist nun nach dem im 31.\ Artikel bewiesenen constant $V^0$ in der ganzen
